\begin{abstract}
本文围绕红外小目标无人机检测组织研究内容，包括数据处理、YOLO26 模型设计、模块对照、检测评价和边缘部署。报告采用统一的数据划分与评价指标组织实验记录，分析浅层特征提取与高分辨率预测对检测效果和计算开销的影响。
\end{abstract}
\begin{IEEEkeywords}
红外无人机检测；YOLO26；小目标；边缘部署
\end{IEEEkeywords}

\section{研究背景与相关工作}
YOLO26 结合端到端检测结构和小目标感知标签分配，为实时框检测提供基础\cite{jocher2026yolo26,ultralytics2026docs}。PConv 从局部采样方式出发增强红外小目标特征\cite{yang2025pconv}。
\IEEEpubidadjcol

\section{数据集与预处理}
主实验数据为 SIRST-UAVB，其无人机和鸟类框用于目标学习与误报分析\cite{yang2025pconv}。
% 填入实际取得的数据版本、来源序列、划分清单与预处理参数。

\section{模型结构与算法设计}
研究以 YOLO26n 为基线，分别考察浅层 PConv 与高分辨率 P2 预测分支。
% 实现完成后补充模块位置、结构图和配置文件索引。

\section{实验设置}
% 填入训练环境、权重初始化、超参数、随机种子和选模规则。

\section{实验结果与分析}
% 填入真实训练所得指标、配对对照、消融、尺度分组和失败案例。

\section{边缘部署与视频验证}
% 填入实际编译版本、量化方式、视频来源、精度变化和端到端延迟。

\section{结论}
% 根据已完成实验总结研究结论。
